\documentclass[conference]{IEEEtran}

%==================== PAQUETES ====================

\usepackage[utf8]{inputenc}
\usepackage[T1]{fontenc}
\usepackage[spanish]{babel}

\usepackage{graphicx}
\usepackage{float}
\usepackage{amsmath}
\usepackage{amssymb}
\usepackage{booktabs}
\usepackage{tabularx}
\usepackage{caption}
\usepackage{cite}
\usepackage{url}

%==================== INICIO ====================

\begin{document}

\title{Título del Documento}

\author{
\IEEEauthorblockN{Carlos Hurtado Villalobos}
\IEEEauthorblockA{
Tecnológico de Costa Rica\\
Ingeniería Electrónica\\
Correo: carlos20032910@gmail.com
}
}

\maketitle

%==================== RESUMEN ====================

\begin{abstract}
Este documento presenta una plantilla base en formato IEEE a doble columna para la elaboración de informes técnicos y académicos.
\end{abstract}

\begin{IEEEkeywords}
IEEE, LaTeX, doble columna, reporte técnico.
\end{IEEEkeywords}

%==================== SECCIONES ====================

\section{Introducción}

Texto de introducción aquí.

\section{Marco Teórico}

Texto del marco teórico.

\subsection{Ejemplo de ecuación}

Ejemplo:

\begin{equation}
V = IR
\end{equation}

\section{Resultados}

%==================== FIGURA ====================

\begin{figure}[H]
\centering
\includegraphics[width=0.8\columnwidth]{imagen.png}
\caption{Ejemplo de figura.}
\label{fig:ejemplo}
\end{figure}

%==================== TABLA ====================

\begin{table}[H]
\centering
\caption{Ejemplo de tabla}
\label{tab:ejemplo}
\begin{tabular}{cc}
\toprule
Dato 1 & Dato 2 \\
\midrule
10 & 20 \\
30 & 40 \\
\bottomrule
\end{tabular}
\end{table}

\section{Conclusiones}

Conclusiones del trabajo.

%==================== REFERENCIAS ====================

\begin{thebibliography}{99}

\bibitem{latex}
L. Lamport,
\textit{\LaTeX: A Document Preparation System}.
Addison-Wesley, 1994.

\bibitem{ieee}
IEEE,
``IEEE Conference Template,''
[Online]. Available:
https://www.ieee.org/

\end{thebibliography}

\end{document}
